% ================== Chapitre 4 : configurations de reference ================
\chapter{Configurations de référence}
\label{ch:configurations}

Les configurations de ce chapitre sont complètes et cohérentes entre elles.
Elles correspondent à la branche 4.9 de PowerDNS Authoritative et à la branche 5
de PowerDNS Recursor. Deux précautions de version s'imposent avant application :
les directives \opt{primary} et \opt{secondary} remplacent les anciennes
directives \opt{master} et \opt{slave} depuis la version 4.5, et la branche 5 du
résolveur introduit un format YAML équivalent au format classique employé ici.
Le format attendu par la version effectivement installée doit être vérifié.

\section{Serveurs autoritatifs ns1 et ns2}
\label{sec:autoritatifs}

\begin{lstlisting}[language=pdnsconf,caption={Configuration du serveur
  autoritatif primaire ns1},label=lst:ns1]
# /etc/powerdns/pdns.conf
# GANDAL T005 : ns1, serveur autoritatif primaire, 10.30.4.10, hote H1

# --- stockage : volume faible, replication assuree par transfert de zone
launch=gsqlite3
gsqlite3-database=/var/lib/powerdns/pdns.sqlite3

# --- ecoute restreinte a l'adresse de service du VNet Services
local-address=10.30.4.10
local-port=53

# --- role
primary=yes
secondary=no

# --- replication : ns2 est le seul destinataire et le seul demandeur autorise
allow-axfr-ips=10.30.4.11/32
only-notify=10.30.4.11
also-notify=10.30.4.11

# --- valeurs par defaut des zones creees (cf. politique de duree de vie)
default-ttl=3600
default-soa-content=ns1.gandal.internal hostmaster.gandal.internal 0 3600 900 604800 300

# --- transport : rester sous la MTU utile de 1450 octets imposee par VXLAN
udp-truncation-threshold=1232
max-tcp-connections=50

# --- discretion et isolement
version-string=anonymous
security-poll-suffix=

# --- API locale ; la terminaison TLS est assuree par le mandataire (4.5)
webserver=yes
webserver-address=127.0.0.1
webserver-port=8081
webserver-allow-from=127.0.0.1/32
api=yes
api-key=$PDNS_API_KEY_NS1

# --- journalisation
loglevel=4
log-dns-queries=no
log-dns-details=yes

# --- execution sous systemd
setuid=pdns
setgid=pdns
daemon=no
guardian=no
\end{lstlisting}

Quatre directives appellent un commentaire. La journalisation des requêtes est
désactivée sur l'autoritatif : ses clients légitimes sont au nombre de quatre et
l'identité réelle du demandeur n'y est pas visible, puisque les requêtes
arrivent par le résolveur. La journalisation utile est placée sur le résolveur,
où l'adresse source est celle du client. La directive \opt{security-poll-suffix}
vidée désactive la vérification périodique de version auprès d'un domaine
externe, requête inutile et indésirable dans un environnement fermé. La
directive \opt{version-string} masque la version du logiciel dans la réponse à
\zone{version.bind}. Enfin, la clé d'API n'est pas écrite dans le fichier : elle
est injectée dans l'environnement du service au démarrage, depuis le magasin de
secrets. À partir de la version 4.7, la directive \opt{api-key} accepte également
une empreinte produite par la commande \opt{pdnsutil hash-password}, ce qui évite
toute forme en clair.

\begin{lstlisting}[language=pdnsconf,caption={Écarts de configuration du serveur
  autoritatif secondaire ns2},label=lst:ns2]
# /etc/powerdns/pdns.conf
# GANDAL T005 : ns2, serveur autoritatif secondaire, 10.30.4.11, hote H3
# Seules les directives qui different de ns1 sont reproduites.

local-address=10.30.4.11

primary=no
secondary=yes

# --- ns1 est la seule source de NOTIFY et de transfert acceptee
allow-notify-from=10.30.4.10/32
xfr-cycle-interval=60
xfr-max-received-mbytes=50

# --- aucun transfert sortant : ns2 ne sert de source a personne
allow-axfr-ips=

api-key=$PDNS_API_KEY_NS2
\end{lstlisting}

La directive \opt{allow-axfr-ips} laissée vide ferme toute possibilité de
transfert depuis ns2. L'API de ns2 reste active pour la supervision et
l'inspection locale, mais le mandataire placé devant elle n'autorise que la
méthode \opt{GET}, de sorte que le propriétaire d'écriture reste unique au sens
de la section~\ref{sec:proprietaire}.

\section{Création des zones et des clés de transfert}
\label{sec:zones}

La clé de transfert est unique parce que l'unique transfert du dispositif va de
ns1 vers ns2. Le document 03 évoque des clés TSIG limitées par zone : cette
granularité devient nécessaire dans la variante où les mises à jour dynamiques
du RFC 2136 seraient employées, chaque émetteur ne devant alors pouvoir modifier
que sa propre zone. Cette variante n'est pas retenue, comme l'explique la
section~\ref{sec:proprietaire}.

\begin{lstlisting}[caption={Création de l'espace de nommage sur ns1},
  label=lst:zones]
# --- sur ns1 : generation de la cle de transfert
pdnsutil generate-tsig-key gandal-xfr hmac-sha256
# la valeur produite est deposee dans le magasin de secrets, jamais dans Git

# --- zone racine interne et identites des serveurs
pdnsutil create-zone gandal.internal ns1.gandal.internal
pdnsutil add-record  gandal.internal @    NS 3600 ns2.gandal.internal.
pdnsutil add-record  gandal.internal ns1  A  3600 10.30.4.10
pdnsutil add-record  gandal.internal ns2  A  3600 10.30.4.11
pdnsutil add-record  gandal.internal rec1 A  3600 10.30.4.12
pdnsutil add-record  gandal.internal rec2 A  3600 10.30.4.13

# --- zone d'infrastructure, puis sa delegation depuis la zone parente
pdnsutil create-zone infra.gandal.internal ns1.gandal.internal
pdnsutil add-record  infra.gandal.internal @ NS 3600 ns2.gandal.internal.
pdnsutil add-record  gandal.internal infra NS 3600 ns1.gandal.internal.
pdnsutil add-record  gandal.internal infra NS 3600 ns2.gandal.internal.

# --- inventaire d'infrastructure, extrait
pdnsutil add-record infra.gandal.internal h1      A 3600 10.30.0.11
pdnsutil add-record infra.gandal.internal h2      A 3600 10.30.0.12
pdnsutil add-record infra.gandal.internal h3      A 3600 10.30.0.13
pdnsutil add-record infra.gandal.internal xo      A 3600 10.30.0.20
pdnsutil add-record infra.gandal.internal bastion A 3600 10.30.0.21
pdnsutil add-record infra.gandal.internal kea1    A 3600 10.30.4.20
pdnsutil add-record infra.gandal.internal kea2    A 3600 10.30.4.21
pdnsutil add-record infra.gandal.internal ntp1    A 3600 10.30.4.30
pdnsutil add-record infra.gandal.internal ntp2    A 3600 10.30.4.31
pdnsutil add-record infra.gandal.internal ipam    A 3600 10.30.4.40
pdnsutil add-record infra.gandal.internal agent   A 3600 10.30.4.41

# --- zones inverses d'infrastructure
for O in 0 1 2 4 5 6 7 8 10 16 17 ; do
  pdnsutil create-zone "${O}.30.10.in-addr.arpa" ns1.gandal.internal
  pdnsutil add-record  "${O}.30.10.in-addr.arpa" @ NS 3600 ns2.gandal.internal.
done

# --- cinquante tenants : zone directe, zone inverse et delegation
for N in $(seq 1 50) ; do
  T=$(printf 'tenant-%03d' "$N")
  O=$(( 63 + N ))
  pdnsutil create-zone "${T}.gandal.internal" ns1.gandal.internal
  pdnsutil add-record  "${T}.gandal.internal" @  NS 3600 ns2.gandal.internal.
  pdnsutil add-record  "${T}.gandal.internal" gw A  3600 "10.30.${O}.1"
  pdnsutil create-zone "${O}.30.10.in-addr.arpa" ns1.gandal.internal
  pdnsutil add-record  "${O}.30.10.in-addr.arpa" @ NS  3600 ns2.gandal.internal.
  pdnsutil add-record  "${O}.30.10.in-addr.arpa" 1 PTR 3600 "gw.${T}.gandal.internal."
  pdnsutil add-record  gandal.internal "$T" NS 3600 ns1.gandal.internal.
  pdnsutil add-record  gandal.internal "$T" NS 3600 ns2.gandal.internal.
done

# --- transfert signe et increment automatique du numero de serie par l'API
for Z in $(pdnsutil list-all-zones) ; do
  pdnsutil set-meta "$Z" TSIG-ALLOW-AXFR gandal-xfr
  pdnsutil set-meta "$Z" SOA-EDIT-API    DEFAULT
done

# --- controle de coherence avant mise en service
pdnsutil rectify-all-zones
pdnsutil check-all-zones
\end{lstlisting}

\begin{lstlisting}[caption={Déclaration des zones secondaires sur ns2},
  label=lst:ns2zones]
# --- sur ns2 : import de la meme cle, puis declaration des zones secondaires
pdnsutil import-tsig-key gandal-xfr hmac-sha256 "$TSIG_GANDAL_XFR"

for Z in $(cat /etc/gandal/zones.list) ; do
  pdnsutil create-secondary-zone "$Z" 10.30.4.10
  pdnsutil activate-tsig-key     "$Z" gandal-xfr secondary
done
\end{lstlisting}

La métadonnée \opt{SOA-EDIT-API} fixée à \opt{DEFAULT} fait incrémenter le
numéro de série à chaque modification reçue par l'API. C'est cette
incrémentation qui déclenche le message \opt{NOTIFY} et le transfert vers ns2 ;
sans elle, les deux serveurs divergeraient silencieusement.

\section{Résolveurs récursifs rec1 et rec2}
\label{sec:recursif}

\begin{lstlisting}[language=pdnsconf,caption={Configuration du résolveur
  récursif rec1},label=lst:rec1]
# /etc/powerdns/recursor.conf
# GANDAL T005 : rec1, resolveur recursif, 10.30.4.12, hote H2
# rec2 est identique, avec local-address=10.30.4.13 sur l'hote H1

local-address=10.30.4.12
local-port=53
threads=2
pdns-distributes-queries=no

# --- clients autorises : strictement les prefixes du plan 02
allow-from=10.30.0.0/24, 10.30.4.0/24, 10.30.5.0/24, 10.30.6.0/24,
           10.30.7.0/24, 10.30.8.0/24, 10.30.16.16/28, 10.30.16.32/28,
           10.30.64.0/18

# --- espace interne : redirection vers les serveurs autoritatifs
forward-zones=gandal.internal=10.30.4.10;10.30.4.11
forward-zones+=30.10.in-addr.arpa=10.30.4.10;10.30.4.11

# Sans cette directive, le resolveur servirait localement les zones inverses
# privees du RFC 6303 et court-circuiterait la redirection ci-dessus.
serve-rfc1918=no

# --- validation DNSSEC de l'espace public, exception traitee dans le Lua
dnssec=validate
lua-config-file=/etc/powerdns/recursor-config.lua

# --- politique de selection des zones (section 4.4)
lua-dns-script=/etc/powerdns/views.lua

# --- transport : rester sous la MTU utile de 1450 octets imposee par VXLAN
edns-outgoing-bufsize=1232
udp-truncation-threshold=1232
qname-minimization=yes

# --- cache dimensionne pour 4 GiB de memoire
max-cache-entries=500000
max-packetcache-entries=250000
max-cache-ttl=86400
max-negative-ttl=300

# --- discretion et isolement
version-string=anonymous
security-poll-suffix=
export-etc-hosts=no

# --- point de metriques, expose par le mandataire (section 4.5)
webserver=yes
webserver-address=127.0.0.1
webserver-port=8082
webserver-allow-from=127.0.0.1/32
api-key=$PDNS_REC_API_KEY

loglevel=4
\end{lstlisting}

La directive \opt{serve-rfc1918} mérite d'être soulignée. Par défaut, le
résolveur se déclare autoritatif et vide pour les zones inverses des espaces
privés, afin d'éviter des requêtes inutiles vers la racine. Dans GANDAL, ce
comportement empêcherait toute résolution inverse interne. Il est donc désactivé
et remplacé par une redirection explicite de \zone{30.10.in-addr.arpa} vers les
autoritatifs.

La directive \opt{export-etc-hosts} est maintenue à la valeur \opt{no} : le
fichier \opt{hosts} local des résolveurs contient le socle de secours décrit à la
section~\ref{sec:amorcage}, qui n'a pas vocation à être publié et qui pourrait
entrer en contradiction avec les zones autoritatives.

\begin{lstlisting}[language=luapdns,caption={Configuration Lua du résolveur :
  exception de validation et journalisation},label=lst:reclua]
-- /etc/powerdns/recursor-config.lua
-- Charge par la directive lua-config-file.

-- Le domaine de premier niveau "internal" est reserve a l'usage prive et n'est
-- pas delegue dans la racine publique. Avec dnssec=validate, le resolveur
-- obtiendrait de la racine une preuve signee de son inexistence et declarerait
-- toute reponse interne invalide. L'exception locale ci-dessous suspend la
-- validation sous cette branche, et sous elle seulement.
addNTA("internal.", "espace de nommage interne GANDAL, non delegue et non signe")

-- Aucune exception n'est necessaire pour 10.in-addr.arpa : cette branche est
-- deleguee sans enregistrement DS dans l'arbre public, donc deja non signee.

-- Journalisation structuree vers un collecteur local du VNet Services.
-- Monitoring lit ensuite les fichiers produits en pull (IF-NS-08).
dnstapFrameStreamServer({"/var/run/pdns-recursor/dnstap.sock"})
\end{lstlisting}

\section{Script de politique de vue}
\label{sec:script}

Le script ci-dessous met en \oe uvre les équations \eqref{eq:tau} et
\eqref{eq:zones}. Il est volontairement court et sans accès externe : il
s'exécute pour chaque requête et toute opération coûteuse s'y traduirait
directement en latence de résolution.

\begin{lstlisting}[language=luapdns,caption={Politique de sélection des zones
  appliquée sur rec1 et rec2},label=lst:views]
-- /etc/powerdns/views.lua
-- Politique de selection des zones par adresse source.
-- Exigences couvertes : NET-FR-008, IF-NS-03 ; scenario de test S04.

local ZONE_RACINE = "gandal.internal"
local SUFFIXE_INV = "30.10.in-addr.arpa"

local TAG_ADMIN = 1     -- gestion, services, monitoring, VPN administrateur
local TAG_REFUS = 999   -- source autorisee mais sans vue interne

local reseauxAdmin = newNMG()
reseauxAdmin:addMask("10.30.0.0/24")
reseauxAdmin:addMask("10.30.4.0/24")
reseauxAdmin:addMask("10.30.5.0/24")
reseauxAdmin:addMask("10.30.16.32/28")

local reseauxTenant = newNMG()
reseauxTenant:addMask("10.30.64.0/18")

-- Plan 02 : le tenant n occupe 10.30.(63+n).0/24, pour 1 <= n <= 50.
local function numeroTenant(adresse)
  local o3 = tonumber(adresse:toString():match("^10%.30%.(%d+)%.%d+$"))
  if o3 == nil or o3 < 64 or o3 > 113 then return nil end
  return o3 - 63
end

-- gettag() est evalue avant toute consultation de cache et sa valeur entre
-- dans la cle du cache paquet : deux vues ne partagent jamais une reponse.
function gettag(remote, ednssubnet, localip, qname, qtype, ednsoptions, tcp)
  if reseauxAdmin:match(remote) then return TAG_ADMIN end
  if reseauxTenant:match(remote) then
    local n = numeroTenant(remote)
    if n ~= nil then return 1000 + n end
  end
  return TAG_REFUS
end

local function sousDomaine(nom, zone)
  return nom == zone or nom:sub(-(#zone + 1)) == ("." .. zone)
end

function preresolve(dq)
  local nom     = dq.qname:toStringNoDot():lower()
  local interne = sousDomaine(nom, ZONE_RACINE)
  local inverse = sousDomaine(nom, SUFFIXE_INV)

  -- Les noms publics suivent la recursion ordinaire, quelle que soit la vue.
  if not interne and not inverse then return false end

  if dq.tag == TAG_ADMIN then return false end

  if dq.tag > 1000 then
    local n           = dq.tag - 1000
    local zoneDirecte = string.format("tenant-%03d.%s", n, ZONE_RACINE)
    local zoneInverse = string.format("%d.%s", n + 63, SUFFIXE_INV)
    if sousDomaine(nom, zoneDirecte) or sousDomaine(nom, zoneInverse) then
      return false
    end
  end

  -- REFUSED plutot que NXDOMAIN : le refus ne revele pas l'existence du nom.
  dq.rcode = pdns.REFUSED
  return true
end
\end{lstlisting}

Le script refuse également l'apex \zone{gandal.internal} aux sources tenant, ce
qui empêche l'énumération des délégations et donc la découverte de la liste des
tenants. Il laisse en revanche passer sans traitement toute requête portant sur
un nom public, afin de ne pas pénaliser le cas le plus fréquent.

\section{Exposition de l'API et des métriques}
\label{sec:api}

Le serveur HTTP intégré à PowerDNS ne fournit pas de terminaison TLS.
L'exigence \exig{net-fr-016}, qui demande des API administratives chiffrées, est
donc satisfaite en liant ce serveur à la boucle locale et en plaçant devant lui
un mandataire qui assure le chiffrement et le contrôle d'accès par adresse.

\begin{lstlisting}[caption={Terminaison TLS et contrôle d'accès devant l'API
  autoritative},label=lst:nginx]
# /etc/nginx/conf.d/pdns-api.conf  sur ns1 (10.30.4.10)

server {
    listen 10.30.4.10:8443 ssl;
    server_name ns1.gandal.internal;

    ssl_certificate     /etc/ssl/gandal/ns1.crt;   # PKI interne, IF-EXT-01
    ssl_certificate_key /etc/ssl/gandal/ns1.key;
    ssl_protocols       TLSv1.2 TLSv1.3;

    # Seul l'agent reseau publie des enregistrements.
    location /api/ {
        allow 10.30.4.41/32;
        deny  all;
        proxy_pass http://127.0.0.1:8081;
    }

    # Seul Monitoring lit les metriques, et il en est l'initiateur.
    location = /metrics {
        allow 10.30.5.10/32;
        deny  all;
        proxy_pass http://127.0.0.1:8081/metrics;
    }

    location / { return 404; }
}

# Sur ns2, le bloc /api/ est complete par la restriction suivante, qui
# garantit l'unicite du proprietaire d'ecriture :
#     limit_except GET { deny all; }
#
# Sur rec1 et rec2, seul le bloc /metrics est publie, vers 127.0.0.1:8082.
\end{lstlisting}

Ce montage répond également à l'interface IF-NS-08 : les quatre instances
n'ouvrent aucune connexion vers Monitoring, qui vient lire leur point de
métriques. La limitation de débit en réponse n'est pas activée sur les serveurs
autoritatifs, et cette absence est délibérée : PowerDNS Authoritative n'intègre
pas cette fonction, et le risque d'amplification qu'elle traite suppose une
exposition publique que la matrice de flux interdit. Si cette exposition devait
changer, un répartiteur DNS spécialisé devrait être placé en amont.
